\documentclass[11pt]{article}
\usepackage{mathunal-base}
\mathunalfuente{serif}

\renewcommand{\examtitulo}{Primer Examen Parcial de Álgebra Lineal}
\renewcommand{\examfecha}{9 de Abril de 2022}

\newcommand{\sino}{\fbox{Sí}\ \fbox{No}}

\begin{document}

\examheader

\vspace{4pt}
\noindent
\begin{minipage}[t]{0.62\linewidth}
\textbf{Puntaje. Sólo para uso Oficial}\\[4pt]
\begin{tabular}{|c|c|c|c|c|c|c|c|c|}
\hline
1 & 2 & 3 & 4 & 5 & 6 & 7 & \textbf{TOTAL} & \textbf{NOTA} \\ \hline
 & & & & & & & & \\ \hline
\end{tabular}
\end{minipage}

\vspace{6pt}
\noindent\textbf{Instrucciones:} La duración del examen es de 1 hora y 50 minutos. El examen consta de siete preguntas en tres hojas impresas por ambos lados, verifique que su examen esté completo. En las preguntas con procedimiento justifique sus respuestas en los espacios asignados. No está permitido sacar hojas en blanco ni ningún tipo de apuntes durante el examen, verifique que su celular esté apagado. No se permite el uso de calculadora.

\vspace{6pt}
\noindent\textbf{IDENTIFICACIÓN}

\vspace{8pt}
\noindent Nombre: \dotfill\ Cédula \dotfill

\vspace{4pt}
\noindent Profesor: \dotfill\ Grupo \dotfill

\vspace{6pt}
\begin{center}
\textbf{Espacio en blanco para realizar cómputos}
\end{center}

\newpage
\noindent\textbf{\large I. Completación}

\vspace{2pt}
\noindent En las preguntas 1 a 6 complete los espacios en blanco. \textbf{NOTA}: En esta sección se califica sólo la respuesta y no se tiene en cuenta el procedimiento.

\begin{enumerate}
\item{}[10pt] Sean $A = \begin{bmatrix}7 & -5\\ 1 & -1\end{bmatrix}$ y $b = \begin{bmatrix}3\\-1\end{bmatrix}$. Encuentre todas las soluciones, si las hay, del sistema de ecuaciones lineales asociado a la matriz aumentada $[A|b]$.

\vspace{4pt}
\fbox{\begin{minipage}[t][2.5cm]{\dimexpr\linewidth-2\fboxsep-2\fboxrule}
$x = $

$y = $
\end{minipage}}

\item{}[10pt] Sean $A = \begin{bmatrix}1 & 0 & -1\\ 1 & 1 & 0\end{bmatrix}$ y $B = \begin{bmatrix}3 & -3\\ 3 & -3\\ 3 & -3\end{bmatrix}$.
\begin{enumerate}
\item{}[5pt] Calcule $AB$.

\vspace{4pt}
\fbox{\begin{minipage}[t][2.3cm]{\dimexpr\linewidth-2\fboxsep-2\fboxrule}
$AB = $
\end{minipage}}

\item{}[5pt] Calcule $B^TA^T$.

\vspace{4pt}
\fbox{\begin{minipage}[t][2.3cm]{\dimexpr\linewidth-2\fboxsep-2\fboxrule}
$B^TA^T = $
\end{minipage}}
\end{enumerate}

\item{}[10pt] Considere la matriz $A = \begin{bmatrix}a & b & c\\ d & e & f\\ g & h & i\end{bmatrix}$ y suponga que $\det(A) = -1$. Calcule los siguientes determinantes.

\vspace{6pt}
\noindent (a) [5pt]

\vspace{4pt}
\fbox{\begin{minipage}[t][3cm]{\dimexpr\linewidth-2\fboxsep-2\fboxrule}
$\begin{vmatrix} g & h & i\\ 2d+5a & 2e+5b & 2f+5c\\ a & b & c\end{vmatrix} = $
\end{minipage}}

\vspace{8pt}
\noindent (b) [5pt]

\vspace{4pt}
\fbox{\begin{minipage}[t][2cm]{\dimexpr\linewidth-2\fboxsep-2\fboxrule}
$\det(2A^3) = $
\end{minipage}}

\item{}[20pt] Lucho vende tinto a \$2000. Suponga que usted solamente tiene monedas de \$50, \$100 y \$200. \textbf{Escriba todas las combinaciones posibles} para pagarle un tinto a Lucho utilizando exactamente trece monedas de las que usted tiene. En otras palabras, describa el número de monedas de cada tipo que se pueden utilizar para pagarle a Lucho con exactamente trece monedas.

\vspace{4pt}
\fbox{\begin{minipage}[t][4cm]{\dimexpr\linewidth-2\fboxsep-2\fboxrule}
\phantom{x}
\end{minipage}}

\item{}[15pt] Considere los vectores $v_1 = \begin{bmatrix}1\\-1\\0\end{bmatrix}$, $v_2 = \begin{bmatrix}2\\-2\\0\end{bmatrix}$, $v_3 = \begin{bmatrix}1\\0\\-1\end{bmatrix}$ y $v_4 = \begin{bmatrix}1\\1\\1\end{bmatrix}$ y sea $A = [v_1|v_2|v_3|v_4]$ la matriz de tamaño $3\times 4$ cuyas columnas son los vectores $v_1,v_2,v_3,v_4$. Después de aplicar operaciones elementales por filas se encuentra que la forma escalonada reducida por renglones de $A$ es $\begin{bmatrix}1 & 2 & 0 & 0\\ 0 & 0 & 1 & 0\\ 0 & 0 & 0 & 1\end{bmatrix}$. Con esta información responda las siguientes preguntas:
\begin{enumerate}
\item{}[5pt] ¿Cuál es el rango de $A$?

\vspace{4pt}
\fbox{\begin{minipage}[t][1.6cm]{\dimexpr\linewidth-2\fboxsep-2\fboxrule}
$\mathrm{Rango}(A) = $
\end{minipage}}

\item{}[5pt] ¿Son los vectores $v_1, v_2, v_3$ y $v_4$ linealmente independientes?

\vspace{4pt}
\sino

\item{}[5pt] Sea $b = \begin{bmatrix}1\\1\\3\end{bmatrix}$. ¿Puede expresarse $b$ como una combinación lineal de $v_1, v_2, v_3$ y $v_4$?

\vspace{4pt}
\sino
\end{enumerate}

\item{}[10pt] Sean $u, v$ vectores en $\mathbb{R}^n$. Suponga que $u$ y $v$ son ortogonales y que $||u||=3$ y $||v||=4$.
\begin{enumerate}
\item{}[5pt] Calcule $\mathrm{Proy}_v(u)$.

\vspace{4pt}
\fbox{\begin{minipage}[t][2cm]{\dimexpr\linewidth-2\fboxsep-2\fboxrule}
$\mathrm{Proy}_v(u) = $
\end{minipage}}

\item{}[5pt] Calcule $||u+v||$.

\vspace{4pt}
\fbox{\begin{minipage}[t][2cm]{\dimexpr\linewidth-2\fboxsep-2\fboxrule}
$||u+v|| = $
\end{minipage}}
\end{enumerate}
\end{enumerate}

\newpage
\noindent\textbf{\large II. Solución con Procedimiento}

\vspace{4pt}
\noindent En esta parte del examen se deben justificar las respuestas.

\begin{enumerate}
\setcounter{enumi}{6}
\item{}[25pt] Considere el sistema homogéneo de ecuaciones lineales dado por:
\[
\left\{
\begin{aligned}
x + y + 2z + w &= 0\\
-x + y + 4z - w &= 0\\
2x + y + z + 2w &= 0
\end{aligned}
\right.
\]
\begin{enumerate}
\item{}[8pt] ¿Es el sistema consistente? En caso afirmativo, ¿Cuántas soluciones tiene el sistema?

\vspace{3.5cm}

\item{}[9pt] Un cálculo en MATLAB muestra que la forma escalonada reducida por renglones del sistema es
\[
\left[\begin{array}{cccc|c}1 & 0 & -1 & 1 & 0\\ 0 & 1 & 3 & 0 & 0\\ 0 & 0 & 0 & 0 & 0\end{array}\right].
\]
Encuentre la solución general al sistema.

\vspace{4cm}

\item{}[8pt] Si es posible escriba dos vectores $v_1$ y $v_2$ que formen un conjunto generador del conjunto solución del sistema. Si no es posible, explique por qué.

\vspace{4cm}
\end{enumerate}
\end{enumerate}

\examfooter
\end{document}
